\id{МНРТИ 06.81.30}{}

\begin{header}
\swa{A.~Amanbekova}{MERGERS AND ACQUISITIONS: THE ROLE OF DUE DILIGENCE IN M\&A OUTCOMES}

A.~Amanbekova\envelope
\end{header}

\begin{affil}
Kazakh-British Technical University, Almaty, Kazakhstan

\corrauthor{Corresponding author: amanbekovaidana@gmail.com}
\end{affil}

The article establishes due diligence as the essential factor which
decides M\&A transaction success through its core purpose. Organizations
need due diligence (hereinafter referred to as ``DD'') as their
essential M\&A success factor because it enables them to reach their
objectives through unstable market conditions and complicated business
operations. The Bain \& Company survey shows that most failed deals
result from various factors which include DD-related issues that lead to
incorrect synergy assessments and failed due diligence that missed
important problems and target companies that received artificial
enhancements for sale purposes and targets that lack proper strategic
alignment and teams ignored their concerns while participating in the
competitive auction process. The research findings show that DD serves
as the fundamental component because it establishes the level at which
organizations achieve their strategic decision implementation regarding
integration and synergy assessment and risk management.

The research uses secondary data analysis to develop multiple proxy
variables which measure DD quality and then evaluates these variables
against their effects on deal achievement. The regression analysis is
performed to assess outcomes by using the change in acquirer stock price
as the dependent variable which measures price changes from
pre-announcement to post-closing periods.

{\bfseries Keywords:} due diligence, mergers and acquisitions, role of due
diligence in deal success, stock price change after deals, market
volatility, regression analysis, share price reaction

\begin{header}
БІРІГУ ЖӘНЕ ЖҰТЫЛУ: БІРІГУ ЖӘНЕ ЖҰТЫЛУ НӘТИЖЕЛЕРІНДЕГІ ДЬЮ ДИЛИДЖЕНСТІҢ РӨЛІ

А.М.~Аманбекова\envelope
\end{header}

\begin{affil}
Қазақстан-Британ Техникалық Университеті, Алматы, Қазақстан,

e-mail: amanbekovaidana@gmail.com
\end{affil}

Бұл мақаланың мақсаты - бірігу және жұтылу мәмілелеріндегі дью дилидженс
процесінің маңызын түсіндіру. Нарықтағы жоғары құбылмалылық және
корпоративтік стратегиялардың барған сайын күрделенуі жағдайында дью
дилидженс (ДД) бірігу және жұтылу (M\&A) мәмілелерінің тиімділігін
қамтамасыз етудегі негізгі кезеңге айналады. Bain \& Company жүргізген
сауалнама нәтижелеріне сәйкес, сәтсіз мәмілелердің көпшілігі бірқатар
себептермен түсіндіріледі, олардың көбі ДД-ге қатысты факторлар болып
табылады: синергияның асыра бағалануы, дью дилидженстің негізгі
мәселелерді анықтай алмауы, мақсатты компанияның сату үшін «әсемделуі»,
стратегиялық сәйкестіктің жеткіліксіздігі, күмәнді белгілердің еленбеуі,
сондай-ақ тендерлік бәсекенің қызуына еріп кету. Осы тұжырымдарға сүйене
отырып, ДД рөлін іргелі деп санауға болады: дәл осы кезеңде интеграция,
синергияны бағалау және тәуекелдерді басқару бойынша кейінгі
стратегиялық шешімдердің сапасы қалыптасады.

Бұл жұмыста екінші реттік деректерді талдау тәсілі қолданылады: ДД
сапасын өлшеу үшін бірнеше прокси-көрсеткіш ұсынылып, олардың мәміленің
табыстылығына әсерін эмпирикалық түрде тексеріледі. Нәтижелерді бағалау
үшін регрессиялық талдау жүргізіледі, мұнда тәуелді айнымалы - мәміле
жарияланғанға дейінгі кезең мен мәміле жабылғаннан кейінгі кезең
арасындағы сатып алушы компанияның акция бағасының өзгерісі.

{\bfseries Түйін сөздер:} дью дилидженс, бірігу және жұтылу, мәміленің
сәтті болуындағы дью дилидженстің рөлі, мәмілелерден кейін акция
бағаларының өзгеруі, нарық құбылмалылығы, регрессиялық талдау, акция
бағаларының реакциясы

\begin{header}
СЛИЯНИЯ И ПОГЛОЩЕНИЯ: РОЛЬ ДЬЮ ДИЛИДЖЕНСА В РЕЗУЛЬТАТАХ СЛИЯНИЙ И ПОГЛОЩЕНИЙ

А.М.~Аманбекова\envelope
\end{header}

\begin{affil}
Казахстанско-Британский Технический Университет, Алматы, Казахстан,

e-mail: amanbekovaidana@gmail.com
\end{affil}

Цель этой статьи - объяснить важность дью дилидженса в сделках слияний и
поглощений. В условиях высокой волатильности рынков и усложнения
корпоративных стратегий дью дилидженс (ДД) становится ключевым этапом
обеспечения эффективности сделок по слияниям и поглощениям (M\&A).
Согласно результатам опроса Bain \& Company, большинство неудачных
сделок объясняется рядом причин, многие из которых являются
ДД-факторами, такими как: переоценка синергий, неспособность дью
дилидженса выявить ключевые проблемы, «приукрашивание» целевой компании
перед продажей, недостаточное стратегическое соответствие, игнорирование
сомнений и предупреждающих сигналов, а также вовлечение в азарт
конкурентных торгов. Исходя из этих выводов, роль дью дилидженса можно
считать фундаментальной: именно на этом этапе формируется качество
дальнейших стратегических решений по интеграции, оценке синергий и
управлению рисками.

В данной статье используется подход анализа вторичных данных,
предлагающий ряд прокси-показателей для измерения качества ДД и
эмпирически оценивающий их влияние на успешность сделки. Для оценки
результатов проводится регрессионный анализ, где зависимой переменной
выступает изменение стоимости акций компании-покупателя между периодом
до объявления сделки и периодом после её закрытия.

{\bfseries Ключевые слова:} дью дилидженс, слияния и поглощения, роль дью
дилидженса в успешности сделок, изменение цены акций после сделок,
волатильность рынка, регрессионный анализ, реакция цены акций

\begin{multicols}{2}
{\bfseries Introduction.} Mergers and acquisitions (here\-inafter referred to
as ``M\&A'') require due diligence (hereinafter referred to as ``DD'')
because M\&A transactions generate unpredictable results which result in
various acquisition deal outcomes. M\&A serve as a strategic growth tool
which appears frequently, and many research show that acquisition
performance for acquiring companies remains uncertain after deal
completion.

For example, a key study by King et al. (2003) {[}1{]} summarizes many
empirical works, which prove that in the average M\&A deals results are
not consistently positive for acquirers, and the variation in outcomes
remains large and partly unexplained. Researchers continue to seek
``moderator factors'' which will help them predict deal success or
failure during the initial preparation phase. Due diligence is important
here, because it is a systematic review of the target company that helps
identify risks, assumptions, and structural issues that can strongly
affect post-deal performance.

Market practice confirms that due diligence is one of the main
``bottlenecks'' in successful deal closing. The KPMG 2025 M\&A Deal
Market Study (survey of 300 deal professionals) {[}2{]} shows that
completing due diligence is among the top three obstacles to closing
deals, in line with valuation and integration-related issues. The study
reports that 41\% of respondents mention completing due diligence as a
critical difficulty during the deal process. This supports the idea that
due diligence is not a formal checklist, but a key stage that can slow
down, complicate, or block a deal. Also, according to this study, the
rising number of deal risks demands advanced valuation models which need
more information and stronger regulatory standards, thus due diligence
serves as a vital assessment tool that affects deal length, market
confidence, and transaction success rates.

Consulting practice also indicates that most M\&A failures result from
integration problems and incorrect synergy assessments, which should be
detected during due diligence evaluation. For example, the Bain report
``The Merger Before the Merger'' {[}3{]} lists key reasons why deals
break down, according to their conducted survey, which include (Fig.1):

- Due diligence failed to highlight key issues

- Target was dressed up for sale

- Insufficient strategic fit

- Doubts were pushed aside

- Caught up in the bidding process

These factors are strongly linked to due diligence quality. They create
``hidden risk'' effects where the buyer overestimates value,
underestimates integration challenges, and closes the deal with
incomplete information.
\end{multicols}

\fig{e/image9}[Fig.1 - Results of the Bain \& Company Survey (n = 250)]

\begin{multicols}{2}
In addition, the McKinsey report "Where mergers go wrong" {[}4{]} also
identifies overestimation of synergies as a stable problem which results
in the "winner' s curse" because the auction winner ends
up paying excessive amounts as a result of unrealistic expectations and
insufficient proof of their assumptions. The process fails to
demonstrate actual project risks and integration challenges and required
resources when due diligence is not fully completed.

Hence, the research findings will have the following practical
significance:

- due diligence factors linked to deal success can transform due
diligence from ``checklist'' to measurable KPIs and predictive risk
markers. If due diligence quality indicators statistically explain
buyer's stock price reaction to the conducted deals, then they can be
translated into management actions as: 1) Due diligence scope and depth
requirements;

- criteria for advisor selection;

- risk adjusted valuation and SPA protection.

The theoretical significance of the research will combine three
literature streams that usually studied separately:

- event studies (market reaction to corporate events);

- accounting signals (post-deal reporting outcomes, such as impairment);

- due diligence literature (due diligence as a tool to reduce
information asymmetry).

Finally, successful M\&A is important not only on the corporate level,
but beyond it. M\&A reallocates capital, restructures industries, and
transforms economies. Research suggests deals shape market structure,
enable exit of weak players, speed up technology renewal, and help
industries adapt to shocks (Andrade, Mitchell \& Stafford, 2001 {[}5{]};
Andrade, 2004 {[}6{]}; Breinlich, 2008 {[}7{]}). Cross-border M\&A is
also part of FDI flows and supports transfer of technology and
management practices (UNCTAD WIR series {[}8{]}). Failures destroy value
and may create political risks and lower trust in cross-border
investments.

In addition, the effectiveness of a due diligence can vary depending on
the mechanism for the purchase price under the Sales and Purchase
Agreement (SPA). In international M\&A practice, the two most common
mechanisms are the ``locked box'' and ``completion accounts''. Due
diligence under the two mechanisms is carried out differently.

Under the locked box mechanism, the purchase price is determined based
on the company's historical financial statements as of a particular date
in the past (the ``locked box date''). This means that the due diligence
exercise will have two phases: 1) a detailed review of the historical
financial statements and the key financial metrics underlying the
purchase price calculation, and 2) a more general review of the period
since the locked box date to ensure that there have been no material
leakage transactions which could allow the sellers to extract value from
the company at the buyer's expense.

On the other hand, under the completion accounts mechanism, the purchase
price is calculated on the basis of the financial performance of target
company as of the date of completion of transaction. Therefore, the
price adjustment is done post-closing and the need of detailed review of
the period between signing and closing of the transaction is reduced.
The target company's historical financial statements and financial
metrics, such as net debt and net working capital, are reviewed and are
further included in the purchase price calculation for the adjustment
for completion accounts.

Thus, the due diligence under locked box mechanism is generally more
detailed, since the purchase price is fixed, and a historical account
has to be verified as to its accuracy and adequacy for the period from
the effective date to the closing price. In a completion accounts
mechanism, the focus is on the historical aspects and the contingent
items which will impact future price adjustments, such as net debt and
net working capital.

{\bfseries Materials and мethods.} This research examines whether higher
due diligence quality leads to improved post-deal performance for the
acquiring company. The central hypothesis proposes that a higher level
of due diligence quality has a positive impact on the acquirer's
subsequent performance.

Research design and stages. The research proceeds through several steps.
The first step was collecting secondary deal and market data, after that
I built DD quality proxies and other control variables. After all data
was collected, I estimated a linear regression model and checked
multicollinearity using VIF. Lastly, I interpreted coefficients and
statistical significance.

Data and materials. For the purposes of this research, I used secondary
data collected from CapitalIQ and Refinitiv. The sample includes 39
deals, and sample size adequacy is justified using MacKinlay (1997)
{[}9{]} power analysis. In MacKinlay's example, event study test power
with sample size around 30 is about 0.54, which has the power to
identify a real effect. Since this study uses 39 observations (larger
than 30), expected statistical power should not be lower than in that
illustration. The research design with this sample size allows
regression analysis through event study methods which follow established
requirements.

Variables. The dependent variable measures post-deal performance for the
acquirer and taken as a stock price change between before deal
announcement and one year after deal closing (long-run returns). This
reflects realization of expectations, ability to achieve synergies,
integration success, and avoidance of overpayment. This approach is
supported by long-run return studies: 1) Loughran \& Vijh (2012)
{[}10{]}: long-run returns of acquirers differ from short-run returns
and depend on deal form and integration dynamics; 2) Rau \& Vermaelen
(1998) {[}11{]}: post-deal performance differs by bidder type and deal
characteristics.

The \emph{independent variables} are DD quality proxies (Table 1). A key
limitation is that direct DD scope and true DD quality are not
observable, as there is no available information about its underlying
scope. Therefore, I used proxy variables:

- \emph{information complexity of the target} - target type (private /
public), industry (whether target is in a different industry than the
buyer), country (cross-border or not);

- \emph{deal duration} - days between announcement and closing. This is
not a direct DD metric, but assumed that more time can allow more
complete DD (based on {[}12{]});

- \emph{deal advisor} - advisor quality is treated as a valid
independent variable because top-tier M\&A advisors can provide better
analysis, better risk detection, and better structuring, affecting both
short-run market reaction and long-run performance. This is supported by
work linking advisor choice to bidder returns (Golubov, Petmezas \&
Travlos {[}13{]}; Rau, 2000 {[}14{]}). The mechanism is reduction of
information asymmetrical and lower risk of overpayment or wrong synergy
assumptions;

- \emph{goodwill impairment} - goodwill impairment in the first 1-3
years after the deal is used as a proxy for DD quality because it
reflects how correct pre-acquisition assessments were, whether expected
synergies were realized, and whether acquirer overpaid for the target
company. The research by Wangerin (2019) {[}12{]} and Potepa \& Thomas
(2022) {[}15{]} shows that poor debt disclosure practices result in
higher impairment costs and lower business performance and wrong fair
value evaluations. Therefore, impairment is treated as an objective
ex-post signal of DD effectiveness.

Lastly, controls include payment method, acquirer size and performance,
and relative deal size.
\end{multicols}

\tcap{Table 1 - Variables for regression}
\begin{longtblr}[
  label = none,
  entry = none,
]{
  width = \linewidth,
  colspec = {Q[196]Q[77]Q[592]Q[225]},
  cells = {c},
  cells = {font = \small},
  row{1} = {font=\bfseries\small},
  hlines,
  vlines,
}
Variable name & Unit & Definition & Source\\
Performance & ln & Natural logarithm of the acquirer's stock price change (pre-announcement vs 1 year after closing) & CapitalIQ, Refinitiv\\
Target type & dummy & Whether the target company is private or public & CapitalIQ, Refinitiv\\
Industry & dummy & Whether the target company from the same industry as acquirer & CapitalIQ, Refinitiv\\
Country & dummy & Is the deal cross-border & CapitalIQ, Refinitiv\\
Deal duration & days & Number of days between deal announcement and deal closing & CapitalIQ, Refinitiv\\
Advisor & dummy & Whether the deal advisor was top-tier (e.g. Morgan Stanley, Goldman Sachs etc.) & CapitalIQ, Refinitiv\\
Impairments & dummy & Wheter the acquirer had a goodwill impairment in the first 1-3 years after the deal closed & CapitalIQ, Refinitiv\\
Payment & dummy & The deal's payment method (cash / equity) & CapitalIQ, Refinitiv\\
Relative deal size & USDm & The deal's size relatively to acquirer's market capitalization & CapitalIQ, Refinitiv\\
Acquirer equity value & USDm & The equity value of acquirer & CapitalIQ, Refinitiv\\
Acquirer ROA & \% & ROA of acquirer & CapitalIQ, Refinitiv\\
Acquirer leverage & \% & Acquirer's debt-to-equity ratio & CapitalIQ, Refinitiv
\end{longtblr}

Regression model. For the purposes of this research I made multifactor
linear regression model:

\begin{equation}
Performance_{i} = \beta_{0} + \beta_{1} \cdot DDQuality_{i} + \beta_{2} \cdot Controls_{i} + \epsilon_{i}
\end{equation}

\begin{multicols}{2}
Where:
1) $Performance_{i}$ = stock price change (pre-announcement vs 1 year after closing);
2) $DDQuality_{i}$ = DD proxies;
3) $Controls_{i}$ = control variables;
4) $\epsilon_{i}$ = error term.

After regression model was runed, I also conducted multicollinearity
test, where VIF = 1.432, which suggests no multicollinearity concerns.

{\bfseries Results and Discussion.} In this study, I evaluate statistical
significance at the 90\% confidence level (α = 0.10). This choice is
based on the event-study econometrics literature, which explains that
return-based tests can have limited statistical power in small samples
and that inference depends on sample size and return volatility.
Following this guidance, I report results that are significant at the
10\% level as economically meaningful signals that should be tested
further in larger samples. (Kothari \& Warner, 2004) {[}16{]}.

Model fit. The model explains about 30\% of the variation in the
dependent variable (R² = 0.302), while the adjusted R² is 0.140 (Table
2). In M\&A return regressions, such values are acceptable because stock
return outcomes are noisy and influenced by many unobservable factors.
Prior evidence suggests that cross-sectional models predicting acquirer
stock-return reactions often have low explanatory power, so these fit
measures are reasonable for this market setting. (Quariguasi Frota Neto
et al., 2025 {[}17{]}).
\end{multicols}

\tcap{Table 2 - Multifactor regression model results}
\begin{longtblr}[
  label = none,
  entry = none,
]{
  cells = {c},
  cells = {font = \small},
  row{1} = {font=\bfseries\small},
  cell{1}{1} = {c=2}{},
  hlines,
  vlines,
}
Regression Statistics & \\
Multiple R & 0.549265181\\
R Square & 0.301692239\\
Adjusted R Square & 0.139509534\\
Standard Error & 0.415057657\\
Observations & 39
\end{longtblr}

\tcap{Table 3 - Multifactor regression model results}
\begin{longtblr}[
  label = none,
  entry = none,
]{
  cells = {c},
  cells = {font = \small},
  row{1} = {font=\bfseries\small},
  hlines,
  vlines,
}
& Coefficients & Standard Error & P-value\\
Target type (dummy) & 0.076671579 & 0.165979733 & 0.647254323\\
Industry (dummy) & 0.031494975 & 0.14875856 & 0.833669167\\
Country (dummy) & -0.097429955 & 0.186519835 & 0.605019448\\
Deal duration & -3.94365E-05 & 0.001044216 & 0.970108425\\
Advisor (dummy) & 0.344301484 & 0.183779515 & 0.070164574\\
Impairments (dummy) & -0.301884692 & 0.154618128 & 0.059679329\\
Payment (dummy) & 0.101490258 & 0.153903257 & 0.5143302
\end{longtblr}

\begin{multicols}{2}
Key DD-related findings (Table 3). Advisor (dummy) has a positive
coefficient (+0.344) and is marginally significant (p ≈ 0.07). The
coefficient on Advisor is +0.344 with p = 0.070. Since p <{}
0.10, this variable is significant at the 90\% confidence level. The
positive sign indicates that working with an experienced (top-tier)
advisor who represents superior deal due diligence and structure leads
to better one-year post-deal stock performance for the acquiring
company.

Impairments (dummy) has a negative coefficient (-0.302) and is also
marginally significant (p ≈ 0.06) at the 90\% confidence level. The
negative sign indicates that deals followed by goodwill impairment
within 1-3 years tend to have worse one-year post-deal stock price
performance. The evidence supports impairment as an ex-post signal which
reveals unmet expectations and possible overpayment and unidentifiable
risks that became apparent after the transaction completion.

Other variables. The regression results show that target type and
industry similarity and country and deal duration and payment method do
not affect the outcome because their p-values are high. So in this
sample and this specification they do not show strong independent
explanatory power.

{\bfseries Conclusion.} The goal of this study is to test whether due
diligence quality affects M\&A outcomes for the acquirer. The study uses
a regression model within an event-study logic, where performance is
measured as the acquirer's stock price change from before the
announcement to one year after deal closing. DD quality is proxied by
deal complexity measures, deal duration, advisor dummy, and post-deal
goodwill impairment. The results demonstrate the following findings. The
analysis shows that the Advisor dummy variable demonstrates a positive
relationship with post-deal performance results at a marginally
significant level (p ≈ 0.07), the Goodwill impairment dummy variable
produces negative results which affect post-deal performance at a
marginally significant level (p ≈ 0.06). Other variables are not
significant. The present sample needs no further control procedures for
the current assessment, as VIF indicates no multicollinearity (VIF ≈
1.43).

The research findings indicate that DD-related conditions create
significance for what happens after a deal is completed. Specifically:
1) Better advisor-related conditions (proxy for stronger analysis and
structuring) are associated with higher post-deal stock price
performance; 2) Goodwill impairment in the first 1-3 years after a deal
is associated with worse post-deal performance and may reflect weak
pre-deal screening, weak DD, or unrealistic synergy/valuation
assumptions.

So, as a result, the research results confirm that due diligence
functions as an invisible factor which links expected deal outcomes to
actual post-acquisition performance.
\end{multicols}

\begin{center}
{\bfseries References}
\end{center}

\begin{refs}
1. King D.R., Dan D.R., Daily C.M., Covin J.G. Meta-analyses of
post-acquisition performance: indications of unidentified moderators//
Strategic Management Journal. -2004. - Vol.25(2). - P.187-200. DOI
10.1002/smj.371.

2. The KPMG 2025 M\&A Deal Market Study. - 2025. URL:
\href{https://kpmg.com/us/en/articles/2025/2025-ma-deal-market-study.html}{https://kpmg.com} .-
Date of access: 03.01.2026

3. Grob J., Meacham M. The merger before the merger/Bain and Company.
-1999. . URL:
\href{https://www.bain.com/contentassets/afa083a686c847d58c7fef8e6d0ef379/bb\_merger\_before\_the\_merger.pdf}{https://www.bain.com} .- Дата обращения:03.01.2026

4. Christofferson S.A., McNish R.S., Sias D.L. Where mergers go
wrong/The McKinsey Quarterly. -- 2004. - № 2.- P.1-7.

5. Andrade G., Mitchell M., Stafford E. New Evidence and Perspectives on
Mergers// Journal of Economic Perspectives. -2001. -Vol.15 (2). -
P.103-120. \href{https://doi.org/10.1257/jep.15.2.103}{DOI
10.1257/jep.15.2.103}.

6. Andrade G., Stafford E. Investigating the economic role of mergers//
Journal of Corporate Finance. - 2004.- Vol.10(1). -P.1-36 DOI
\href{https://doi.org/10.1016/S0929-1199(02)00023-8}{10.1016/S0929-1199(02)00023-8}.

7. Breinlich H. Trade liberalization and industrial restructuring
through mergers and acquisitions// Journal of International Economics.
-2008. -Vol.76(2). - P.254-266. DOI
\href{https://doi.org/10.1016/j.jinteco.2008.07.007}{10.1016/j.jinteco.2008.07.007}.

8. UNCTAD World Investment Report. International investment in the
digital economy. -- 2025.-262 p. ISBN 978-92-1-003558-3.

9. MacKinlay A.C. Event Studies in Economics and Finance// Journal of
Economic Literature. -1997.-Vol.35(1). - P.13-39.

10. Moeller S.B., Schlingemann F.P., Stulz R.M. Firm size and the gains
from acquisitions// Journal of Financial Economics. -2004. -Vol.73(2).
-P.201--228. DOI
\href{https://doi.org/10.1016/j.jfineco.2003.07.002}{10.1016/j.jfineco.2003.07.002}.

11. Rau P.R., Vermaelen T. Glamour, value and the post-acquisition
performance of acquiring firms// Journal of Financial Economics. -1998.
-Vol.49(2). - P.223--253. DOI
\href{https://doi.org/10.1016/S0304-405X(98)00023-3}{10.1016/S0304-405X(98)00023-3}.

12. Wangerin D. M\&A Due Diligence, Post-Acquisition Performance, and
Financial Reporting for Business Combinations// Contemporary Accounting
Research. -2019.-Vol.36(4). - P.2344--2378. DOI
\href{https://doi.org/10.1111/1911-3846.12520}{10.1111/1911-3846.12520}.

13. Golubov A., Petmezas D., Travlos N.G. When It Pays to Pay Your
Investment Banker: New Evidence on the Role of Financial Advisors in
M\&As// The Journal of Finance. -2012. -Vol.67(1). P.271--311. DOI
\href{https://doi.org/10.1111/j.1540-6261.2011.01712.x}{10.1111/j.1540-6261.2011.01712.x}.

14. Rau P.R. Investment bank market share, contingent fee payments, and
the performance of acquiring firms// Journal of Financial Economics.
-2000. -Vol.56(2). - P.293--324.
DOI:\href{https://doi.org/10.1016/S0304-405X(00)00042-8}{10.1016/S0304-405X(00)00042-8}.

15. Potepa J., Thomas J. Goodwill Impairment After M\&A:
Acquisition-Level Evidence// SSRN Electronic Journal. -2022. -Vol.8(2).
-P.131-155. DOI
\href{http://doi.org/10.2139/ssrn.4239812}{10.2139/ssrn.4239812}.

16. Kothari S.P., Warner J.B. The Econometrics of Event Studies// SSRN
Electronic Journal. -2004. --Vol.1. DOI
\href{https://doi.org/10.2139/ssrn.608601}{10.2139/ssrn.608601}.

17. Neto J.Q.F., Bozos K., Dutordoir M., Nikolopoulos K. Are acquirer
stock price reactions to M\&A announcements in any way predictable? A
machine-learning analysis// Journal of the Operational Research Society.
-2025. -P.1-23. DOI
\href{https://doi.org/10.1080/01605682.2025.2562956}{10.1080/01605682.2025.2562956}.
\end{refs}

\begin{info}
\hspace{1em}\emph{{\bfseries Information about the author}}

Amanbekova A. - Master's student, Kazakh-British Technical University,
Almaty, Kazakhstan, e-mail: amanbekovaidana@gmail.com.

\hspace{1em}\emph{{\bfseries Информация об авторе}}

Аманбекова А. М. - Магистрант, Казахстанско-Британский Технический
Университет, Алматы, Казахстан, e-mail: amanbekovaidana@gmail.com.
\end{info}

